\documentclass[11pt]{article}
\usepackage[margin=25mm]{geometry}
\usepackage{fontspec,unicode-math}
\setmainfont{Latin Modern Roman}
\setmathfont{Latin Modern Math}
\usepackage{amsmath,amssymb,mathtools}
\usepackage{xurl,hyperref}
\hypersetup{colorlinks=true,urlcolor=blue,linkcolor=blue}
\setcounter{secnumdepth}{0}
\setlength{\parindent}{0pt}
\setlength{\parskip}{5pt}
\setlength{\emergencystretch}{3em}
\newcommand{\sref}[2]{\hyperlink{source:#1:#2}{[S#2]}}
\newcommand{\cataloguescope}[1]{\paragraph{Source question.}#1}
\newcommand{\floor}[1]{\left\lfloor#1\right\rfloor}
\begin{document}
\section{Security of Interactive Lattice-Based Group Key Exchange}\label{statement}

\subsection{Definitions and mathematical statement}

\paragraph{Shared notation.}Standard number systems are denoted by $\mathbb N,\mathbb Z,\mathbb Q,\mathbb R,\mathbb C$; index ranges, computational models and specialized notation follow the cited problem source and the definitions below.


A lattice generated by independent vectors $b_1,\ldots,b_d$ is $L=\{\sum_i z_i b_i:z_i\in\mathbb Z\}$. SVP seeks a shortest nonzero lattice vector; SIVP seeks $d$ independent short vectors, with approximation measured against the relevant successive minimum.

Learning with errors gives noisy linear equations $b_i=\langle a_i,s\rangle+e_i\pmod q$ for a hidden vector $s$. Ring-LWE uses the analogous equations in a polynomial quotient ring, under specified sample and error distributions.

A cryptographic security requirement compares an efficient adversary\textquotesingle s advantage with a negligible function of the security parameter. A negligible function is smaller than every inverse polynomial eventually; computational and information-theoretic security are different requirements.
\paragraph{Statement recorded in the supplied source.}
Can provable security be established for the interactive multiparty key exchange protocol of Ding, Xie, and Lin based on LWE and RLWE, viewed as a generalization of their two-party protocol? \sref{statement}{2}
\paragraph{Scope qualification.}
The full protocol and security experiment are not reproduced in the supplied question; the report\textquotesingle s information and adversary models are required.
\paragraph{Source status.}
The supplied catalogue labels this record open. Its source-status metadata was checked against the supplied snapshot on 2026-10-11. This import does not independently establish current open status or mathematical correctness. The catalogue literature assessment is dated 2026-08-21.
\subsection{Short English statement}
Prove security of the cited interactive multiparty lattice key-exchange protocol.
\subsection{Sources}
\begin{itemize}
\item[\hypertarget{source:statement:1}{S1}] Ulam AI and the UnsolvedMath Contributors, UnsolvedMath: A Curated Collection of Open Mathematics Problems, catalogue entry. Catalogue curation and metadata: CC BY 4.0. The supplied source question is retained; definitions and the English summary are editorial additions.
\url{https://huggingface.co/datasets/ulamai/UnsolvedMath}.
\item[\hypertarget{source:statement:2}{S2}] Cryptography (2018), Oberwolfach Reports, DOI 10.4171/owr/2017/5, pp. 213--266
\url{https://doi.org/10.4171/owr/2017/5}.
\item[\hypertarget{source:statement:3}{S3}] Ding et al., A Simple Provably Secure Key Exchange
\url{https://archive.ymsc.tsinghua.edu.cn/pacm_download/672/12694-dingjt-i28.pdf}.
\item[\hypertarget{source:statement:4}{S4}] Oberwolfach Report 5/2017
\url{https://ems.press/doi/pdf/10.4171/OWR/2017/5}.
\textit{Underlying Oberwolfach Reports material retains its source-specific licence, including CC BY-SA 4.0 where applicable; see THIRD\_PARTY\_NOTICES.md.}
\end{itemize}
\paragraph{Preparation and review.}
Original question attributed to the cited source and UnsolvedMath contributors. Codex (GPT-6; exact serving revision not exposed) assisted with definitions, English summaries, source inspection and formatting. No new human mathematical authorship or independent proof verification is asserted. Catalogue identity checked against the complete supplied registry and existing repository statements; known flagged repetitions and obvious identical conjectures excluded. Full mathematical-equivalence/literature review has not been performed. Statement preparation and final review followed contribution prompts C and D; automated checks verify structure and identity, not open status.
\label{end:statement}
\end{document}
